\documentclass[DIN, pagenumber=false, fontsize=11pt, parskip=half]{scrartcl}

\usepackage{amssymb}
\usepackage[spanish]{babel}
\usepackage[utf8]{inputenc}
\usepackage[T1]{fontenc}
\usepackage{textcomp}

% for code
\usepackage{listings}
\usepackage{listingsutf8}
\usepackage{xcolor}

\definecolor{green}{RGB}{64,160,43}
\definecolor{gray}{RGB}{76,79,105}
\definecolor{purple}{RGB}{136, 57, 239}
\definecolor{background}{RGB}{239,241,245}

\lstdefinestyle{mystyle}{
	backgroundcolor=\color{background},
    commentstyle=\color{gray},
    keywordstyle=\color{purple},
    numberstyle=\tiny\color{gray},
    stringstyle=\color{green},
    basicstyle=\sffamily\scriptsize,
    breakatwhitespace=false,
    breaklines=true,
    captionpos=b,
    keepspaces=true,
    numbersep=5pt,
    showspaces=false,
    showstringspaces=false,
    showtabs=false,
    tabsize=2,
	frame=single
}

\lstset{style=mystyle}
\lstset{inputencoding=utf8/latin1}

\setlength{\parindent}{0em}

% set section in CM
\setkomafont{section}{\normalfont\bfseries\Large}

\newcommand{\mytitle}[1]{{\noindent\Large\textbf{#1}}}

%===================================
\begin{document}

\noindent\textbf{Análisis Numerico} \hfill \textbf{UNRC}\\
2242 \hfill Hilario Yáñez Bertola\\

\mytitle{Práctica 1: Ecuaciones en una variable \hfill \small{\today}}

\begin{enumerate}

    \item Para encontrar una solución positiva a $\tan (x) = 2x$, consideremos $f(x) = \tan (x) - 2x$.
    Es claro que $f(\frac{1}{2})<0$ ya que $\tan (\frac{1}{2}) < \tan (\frac{\pi}{4}) = 1 = 2 \cdot 1$,
    y se puede ver que $f(\frac{3}{2}) > 0$. Como $\tan (x)$ es continua en $\mathbb{R} \setminus \bigl\{ \pi \cdot (k + \frac{1}{2}), k \in \mathbb{Z} \bigr\}$,
    también lo será $f$ gracias a la continuidad de $2x$, y lo será en particular en $[\frac{1}{2}, \frac{3}{2}]$.
    De esta manera, estamos en condiciones de aplicar el método de bisección para aproximar una raíz de $f$.
    Luego, para que el error relativo sea menor que $10^{-5}$ tenemos que:

    \[ \frac{|e_n|}{|p|} < \frac{\frac{\frac{3}{2} - \frac{1}{2}}{2^n}}{\frac{3}{2}}
    = \frac{1}{3 \cdot 2^{n-1}} < 10^{-5} \iff \log_2 (\frac{10^5}{3}) + 1 < n \]

    Usando que $\lceil \log_2(\frac{10^5}{3}) \rceil = 16$, se deduce que para $n > 17$ se
    obtiene una aproximación tan chica como se quiere.

    Con la implementación adjunta al final del documento, se obtiene el resultado $x \approx 1.16556$.

    \lstinputlisting[language=Python, firstline=4, lastline=26]{single-variable.py}

\end{enumerate}

\end{document}